\documentclass[../../../include/open-logic-section]{subfiles}

\begin{document}

\olfileid{his}{set}{mythology}

\olsection{ইতিহাসের চেয়ে পুরাণকথা বেশি?}

এক শতাব্দীরও বেশি দূরত্ব থেকে
এই ঘটনাগুলির দিকে ফিরে
তাকালে এই মূল্যায়নগুলি
বিনা প্রশ্নে মেনে নেওয়া
ঠিক হবে না। ফলগুলি
নিশ্চয়ই নতুন, রোমাঞ্চকর
ও বিস্ময়কর ছিল। কিন্তু
সেগুলি আসলে কতটা
আঘাত দিয়েছিল? আর
জ্যামিতিক স্বজ্ঞার উপর
নির্ভর করা উচিত নয়—এটি
কি সত্যিই সেগুলি দেখিয়েছিল?

বিস্ময়ের প্রশ্নে
\citet{Gouvea2011}
দেখিয়েছেন, ডেডেকিন্ডকে
কান্টরের বিখ্যাত
“\emph{je le vois,
mais je ne le crois
pas}” কথাটি
অনেকটা প্রসঙ্গ
থেকে বিচ্ছিন্ন
করে উদ্ধৃত হয়।
আরও কিছু প্রসঙ্গ
দেখা যাক
(\citeauthor{Gouvea2011}
থেকে উদ্ধৃত):
\begin{quote}
বিষয়টি নিয়ে আমার
উৎসাহের জন্য আমাকে
ক্ষমা করবেন, যদি
আপনার সদয়তা ও
ধৈর্যের কাছে এত
বেশি দাবি করি;
সম্প্রতি আপনাকে
যা জানিয়েছি,
তা আমার নিজের
কাছেও এত
অপ্রত্যাশিত ও
নতুন যে, সম্মানিত
বন্ধু, সেগুলির
শুদ্ধতা নিয়ে
আপনার রায় না
পাওয়া পর্যন্ত
মনে শান্তি
পাচ্ছি না।
আপনি আমার
সঙ্গে একমত
না হওয়া পর্যন্ত
আমি শুধু বলতে
পারি:
\emph{je le vois,
mais je ne le
crois pas.}
\end{quote}
কান্টর জানতেন,
তাঁর ফলটি
“এত অপ্রত্যাশিত,
এত নতুন”।
কিন্তু ফলটিকে
তিনি কখনও
\emph{অবিশ্বাস্য}
মনে করতেন
কি না, তা
সন্দেহজনক।
\citeauthor{Gouvea2011}
যেমন দেখিয়েছেন,
তিনি শুধু নিজের
দেওয়া প্রমাণটি
পরীক্ষা করে
দেখতে
ডেডেকিন্ডকে
অনুরোধ
করছিলেন।

জ্যামিতিক
স্বজ্ঞার
প্রশ্নে:
পেয়ানো তাঁর
স্থানপূরক
বক্ররেখার
প্রবন্ধে
কোনো ছবি
দেননি।
কিন্তু
হিলবার্ট
তাঁর
বক্ররেখা
প্রকাশের
সময়ে
নিজের
উদ্দেশ্য
বোঝান:
পাঠকেরা
“নিচের
জ্যামিতিক
স্বজ্ঞার
সাহায্য
নিলে”
পেয়ানোর
ফলটি
স্পষ্টভাবে
বুঝতে
পারবেন;
তারপর
তিনি
\olref[his][set][pathology]{sec}-তে
দেওয়া
ছবিগুলির
মতো
কয়েকটি
\emph{চিত্র}
দেন।

আরও
সাধারণভাবে:
প্রকাশিত
প্রমাণে
চিত্রের
প্রচলন
কিছুটা
কমলেও,
গণিতবিদেরা
প্রমাণের
সময়ে
\emph{প্রায়ই}
চিত্র
ব্যবহার
করেন—এটি
অস্বীকার
করা
যায় না।
(মোটামুটি
বললে:
ভালো
গণিতবিদেরা
জানেন,
কখন
জ্যামিতিক
স্বজ্ঞার
উপর
নির্ভর
করা
যায়।)

সুতরাং
প্রচারিত
কথা
বিনা
বিচারে
বিশ্বাস
কোরো না;
অন্তত
শুধু
শুনেই
মেনে
নিয়ো
না।
আরও
জানতে
\citet{Giaquinto2007}
পড়তে
পারো।

\end{document}
